\scp{Shared innovations}{07-03}
All languages of central Maluku,
except \tsc{Taliabu}, share a merger of PMP *a/*-ay > \it{a}.
This merger is so widespread that \citet[264]{Blust1993} sees
glide truncation (the loss of the coda in historical vowel-glide
sequences) as “one of the most distinctive characteristics of
the phonological history of CMP languages”,
and consequently it has very little value for subgrouping within “CMP”.
In \srf{13-01-03} we also show this merger does not define “CMP”.
However, while *a/*-ay > \it{a} is also found widely in individual
nodes or languages in the \tsc{\ili{Bima}-Lembata},
\tsc{Timor-Babar}, \tsc{Aru}, and \tsc{Seram-Tanimbar-Bomberai}
subgroups, it is worth noting that this particular merger affects
all nodes of \tsc{Sula-Buru}, \tsc{\ili{Ambon}-Seram}, and \tsc{Seram-Tanimbar-Bomberai}.

Similarly, nearly all languages have *-aw > \it{a},
but several languages of the \tsc{Northwest Seram} node of \tsc{\ili{Ambon}-Seram}
retain *aw unchanged or have *aw > \it{akw} (see \srf{10-02}),
preventing us from assigning *-aw > \it{a} to Proto-\ili{Ambon}-Seram.
\tsc{Taliabu} is different from other languages of central Maluku
in having *-ay > \it{e}
and *-aw > \it{o} -- both of which are characteristic of \tsc{Celebic}
languages to the west.
This is discussed further in \crf{ch:Tal}.

While \ili{Banda} does have *-ay/*-aw > \it{a},
this outcome does not fully merge with reflexes of *a in final
syllables, as *a > \it{a {\tl} o}) in final syllables.
The change \mbox{*-ay/*-aw} > \it{a} must have occurred independently
of the same change in other subgroups of the region.
That is, *-ay/*-aw > \it{a} in \ili{Banda} is due to parallel development.
This parallel development may have been due to contact with other
subgroups in the region,
but *-ay/*-aw > \it{a} cannot have occurred once in a common
ancestor of \ili{Banda} and any of the other subgroups in the region.
This is part of the basis for treating \ili{Banda} as an isolate.
See \srf{12-05} for details.

Languages in central Maluku also share a merger of PMP *n/*ñ > \it{n}.
Again, this merger is so common that it has very little value for subgrouping.
Data that clearly reflect *ñ are sparse,
and there are lexically specific traces of *ñ with raised vowels
in some words in some languages in both the \tsc{Sula-Buru}
and \tsc{\ili{Ambon}-Seram} subgroups, indicating that the distinction
must be reconstructed for some etyma, and the merger probably
cannot be assigned to a common ancestor (see \srf{13-01-01}).

